% P16-C4: zero-support semantics, licensed extraction and scalar evaluation.
\section{Identification when observed masses can be zero}
\label{sec:p16-identification-with-zero-support}

The 2006 identification paper permits nonpositive observed laws, requires
positive conditioning events, and separately writes an action-support
condition \cite[p.~1220 and Lemma~1, p.~1221]{ShpitserPearl2006JointIdentification}.
These requirements must be attached to the actual assignments used by
a derivation. This section gives a finite sound extension of the two
identification procedures, including zero prefixes and zero summands.
It also gives an explicit counterexample to the stronger assertion that
support of the original action alone licenses every returned graph formula.

\begin{hbdefinition}{def:p16-zero-support-data}{Models, queries and two quantified support conditions}
Fix a finite graph \(G\), finite nonempty observed value sets
\(\mathbf S=(S_i)_{i\in N}\), and a normalized nonnegative law
\(p\) compatible with the independent-latent-input model class of
\cref{def:p16-latent-input-model}, with the intervention kernels of
\cref{def:p16-latent-q-kernel}.
The compatible models compared below have this same graph, alphabet
family and observed law; zero cells are allowed. For disjoint
\(X,Y\), every \(x\in S_X\) defines an intervention by
equation replacement. A mass \(P_x(y)\) is identified at this
\(p\) if it agrees in every compatible model.

Put \(B_X=\operatorname{pa}_G(X)\setminus X\), and define
\[
 \begin{split}
 \mathcal A_{\mathrm{obs}}(p,X)
   &=\{x\in S_X:\ p(x,b)>0\text{ for every }b\in S_{B_X}
                                     \text{ with }p(b)>0\},\\
 \mathcal A_{\mathrm{all}}(p,X)
   &=\{x\in S_X:\ p(x,b)>0\text{ for every }b\in S_{B_X}\}.
 \end{split}
\]
The first interprets positive action probability on every observed
parent stratum. The second requires a defined positive conditional
on every formal parent assignment. The source writes
\(P(x\mid\operatorname{pa}_G(X)\setminus X)>0\) without
expanding this assignment quantifier. Both readings are displayed;
neither changes the structural definition of an intervention.
Its statement about every action considered may also include actions
introduced inside a derivation, whose support must then be checked there.
\end{hbdefinition}

\begin{hbtheorem}{thm:p16-zero-support-quantifiers}{proved=annex:p16.zero-support-quantifiers}{Which universal support assumptions imply positivity}
Requiring every value of every joint action \(I\subseteq N\)
to belong to \(\mathcal A_{\mathrm{obs}}(p,I)\) is equivalent
to \(p(v)>0\) at every full assignment. This quantifier includes
the full action \(I=N\). The same is true with
\(\mathcal A_{\mathrm{all}}\). Requiring all values of singleton
actions alone, or even of every proper action set, need not imply
full positivity in the latent-input class.
Nor does requiring all values of just the original action set.

For example, \(A=B=U\) with \(U\) a fair bit and graph
\(A\leftrightarrow B\) satisfies both singleton conditions
for both values, but its off-diagonal observed masses are zero.
\end{hbtheorem}

\begin{hbtheorem}{thm:p16-zero-front-door-obstruction}{proved=annex:p16.zero-front-door-obstruction}{Original-action support does not license missing front-door strata}
In graph \(A\to M\to Y,A\leftrightarrow Y\), let
\(U,E\) be independent bits with \(U\) fair and
\(\mathbb P(E=1)=1/4\). Both models use \(A=U,M=A\).
Model 1 uses \(Y=\mathbf1_{\{M\oplus U=1\text{ or }E=1\}}\);
model 2 uses \(Y=E\). Their common observed law is
\[
 p(a,m,y)=
 \begin{cases}
  3/8,&m=a, y=0,\\
  1/8,&m=a, y=1,\\
  0,&m\ne a.
 \end{cases}
\]
Both action values \(a\) satisfy both quantified original-action
support conditions. Nevertheless
\[
 P^1_a(Y=1)=5/8,\qquad P^2_a(Y=1)=1/4.
\]
The ordinary graph trace has no hedge and returns the front-door
expression. Its unsupported conditional \(p(Y=1\mid A=1-a,M=a)\)
has an arbitrary version \(t\in[0,1]\); substitution makes that
expression \(1/8+t/2\). This variation is not removed by summing
only observed positive cells: the missing conditional has weight \(1/2\)
in the required intervention calculation.
\end{hbtheorem}

\subsection{Conditional versions and justified district extraction}

\begin{hbdefinition}{def:p16-zero-conditional-versions}{Versions of prefix conditionals}
Let \(q_R\) be a known, normalized nonnegative current kernel
on random vertices \(R\), with its outside intervention values fixed.
It is the actual law of the retained structural model on \(G[R]\).
Use a topological order \(r_1,\ldots,r_k\) and prefixes
\(R_j=\{r_1,\ldots,r_j\}\). Put
\(\rho_j=\Sigma_{R\setminus R_j}q_R\), \(\rho_0=1\).
A family of \emph{conditional versions} consists of normalized
nonnegative tables \(\theta_j(v_{r_j}\mid v_{R_{j-1}})\)
satisfying
\[
 \rho_j(v)=\rho_{j-1}(v)\theta_j(v_{r_j}\mid v_{R_{j-1}}).
\]
At a positive prefix this fixes the ordinary ratio. At a zero
prefix all its next numerators are zero, and any normalized next
table is a version. For a maximal district \(D\) of \(G[R]\)
write \(h_D^\theta(v)=\prod_{j:r_j\in D}\theta_j(v)\).
Outside arguments, suppressed here, remain fixed parameters in every table.
\end{hbdefinition}

\begin{hbtheorem}{thm:p16-zero-version-factorization}{proved=annex:p16.zero-version-factorization}{All-zero-prefix chain identities and model-compatible versions}
Conditional versions exist, and every choice satisfies
\(q_R=\prod_{j=1}^k\theta_j=\prod_D h_D^\theta\) at every
assignment, including zero cells. Each \(h_D^\theta\) is normalized
on \(S_D\) at fixed outside values. For every structural model
whose current law is \(q_R\), there is a choice of these versions
for which \(h_D^\theta=Q[D]\) simultaneously for all districts
and all assignments. Arbitrary versions need not have this structural
meaning outside observed support.

At a fixed full assignment let \(j_0\) be the first index with
\(\rho_{j_0}=0\), if there is one. The following finite rule
decides whether \(h_D^\theta(v)\) is independent of all versions:
multiply the fixed ratios at the indices of \(D\) before or at
\(j_0\); if their product is zero, the answer is zero. Otherwise
the value is fixed exactly when every later index in \(D\) has
a singleton value set. If there is no zero prefix, all ratios are
fixed. In particular, a first zero at a vertex of \(D\) certifies
\(Q[D](v)=0\); a zero first created by another district does
not alone justify that conclusion.
\end{hbtheorem}

\begin{hbdefinition}{def:p16-zero-extraction-overlap}{A complementary-action overlap check}
For a maximal district \(D\subseteq R\), put \(J=R\setminus D\)
and \(B=\operatorname{pa}_{G[R]}(J)\setminus J\subseteq D\).
For the particular outside assignment \(j\in S_J\) needed by
an extraction, define
\[
 \mathcal O(q_R,D,j):\qquad
 q_R(j,b)>0\quad\text{for every }b\in S_B\text{ with }q_R(b)>0.
\]
The other coordinates of \(D\) are summed in this marginal.
When \(J\) is empty this condition is \(1>0\).
It is the observed-parent-stratum reading of the source-style action
condition for this complementary action in the \emph{current} model.
The stronger reading requiring every formal parent tuple also implies it.
It concerns the requested value \(j\), not every intervention value.
\end{hbdefinition}

\begin{hbtheorem}{thm:p16-zero-masked-extraction}{proved=annex:p16.zero-masked-extraction}{Finite district recovery with exact support licenses}
If \(\mathcal O(q_R,D,j)\) holds, then the actual district
kernel on this parameter slice is
\[
 Q[D](v_D;j)=
 \begin{cases}
  h_D^\theta(v_D,j),&q_R(v_D,j)>0,\\
  0,&q_R(v_D,j)=0.
 \end{cases}
\]
The proof includes a finite prefix-repair argument: with the district's
latent inputs fixed, positive-parent-stratum overlap permits successive
choices of complementary latent inputs that realize the requested action.
Thus a district outcome possible after that action has positive observed
mass together with it, even when some other parent tuples are null.
All factors in the first case have positive denominators, and the
second case contains no conditional ratio. The result is independent
of all arbitrary conditional versions.

More generally, define \(\mathsf{Recover}(q_R,D,j)\) as follows.
At each \(v_D\), first use the version-invariance rule of the
previous theorem when it determines the cell. If any cells remain,
the overlap check, when true, fills their zero-support values by zero.
If every cell is now determined, return the full recovered table;
otherwise return a support record listing the undetermined cells,
their first zero prefixes, the sets \(R,D,J\), the requested
\(j\), the fixed outside parameters, and the failed overlap strata.
Every returned table is exactly \(Q[D]\), normalized and
nonnegative. This is a finite operation on the known current table;
no partial table is silently used as an input kernel.
\end{hbtheorem}

\begin{hbtheorem}{thm:p16-zero-certain-action}{proved=annex:p16.zero-certain-action}{An action already certain in the current model}
If \(q_R(X=x)=1\), then for every response assignment
\(P^R_x(y)=q_R(y)\). This remains valid with hidden common
inputs and zero observed masses. It can settle a specified query
even when the unchanged graph procedure would return a hedge.
\end{hbtheorem}

\subsection{Finite support-aware execution}

\begin{hbdefinition}{def:p16-zero-evaluation-results}{Resolved masses and support records}
A scalar request specifies a whole known current kernel \(q_R\)
at its fixed outside assignment, current action values \(x\), and
response values \(y\). Its result is either \(\mathsf{Known}(a)\),
meaning the exact mass \(a\) in every compatible current model,
or \(\mathsf{Unresolved}\) with a support or graph record.
An unresolved record is not a probability or a proof of nonidentification
at this particular observed law.
Records retain the original query and the containing scalar call frames,
including each random/fixed set and its assigned values.

For a finite product of required causal masses, a proved-zero factor
certifies the product as zero, irrespective of other unresolved factors.
Otherwise all factors must be known before multiplication. A finite
sum is returned only when every summand is thus determined. A recorded
failure is never itself multiplied or divided. All factors represent
actual finite probabilities, so the zero-product inference is legitimate.
\end{hbdefinition}

\begin{hbmethod}{def:p16-zero-support-algorithms}{Scalar extensions of ID and the modified Tian procedures}
Use \cref{def:p16-id-algorithm,def:p16-tian-component-algorithm,def:p16-tian-outer-algorithm}
with the following explicit evaluation rules. Retain the current
random/fixed sets and actual parameter values.
\begin{enumerate}
\item Before an ID scalar request, if \(q_R(X=x)=1\), return
      \(\mathsf{Known}(q_R(y))\). Empty action sets are included.
\item ID ancestral reduction uses the full known marginal table.
      Additional nonancestor actions may use any reference tuple;
      one may try all such tuples and use any known return. If none
      succeeds, retain the unresolved records. Each choice is finite.
\item ID district splitting evaluates every required summand and its
      component scalar factors with the original assigned coordinates.
      Apply the proved-zero product rule before requiring all remaining
      factors. A graph failure becomes an unresolved graph record.
\item In ID's maximal-district base and larger-district recursion,
      replace prefix-ratio extraction by \(\mathsf{Recover}\),
      at the needed values on the current complementary set.
      A recovered full table permits the usual marginal or recursive
      call. A support record makes that scalar request unresolved.
\item For \(\mathsf{c\mbox{-}identify}(C,T,q_T)\), its ancestral
      base returns the requested cell of the known marginal; its
      graph failure is unresolved. Its recursive step forms the full
      ancestral marginal \(q_A\), calls
      \(\mathsf{Recover}(q_A,T',v_{A\setminus T'})\), and
      recurses only if that full table is recovered. One may also
      first apply the certain-action identity to \(T\setminus C\).
\item The outer Tian procedure evaluates its finite sum over
      \(D\setminus Y\) and its district factors. For each factor,
      recover the required full-graph district table at its actual
      complementary assignment, then use the scalar component call.
      Use the same proved-zero product rule. Values outside \(D\cup X\)
      may be fixed arbitrarily; all finitely many choices may be tried.
\end{enumerate}
A requested whole response kernel at \(x\) is returned only when
all its finitely many response masses are known. Every new kernel
passed to recursion is a complete normalized table on its random set
at its specified fixed parameters. The procedure never fills an
undetermined cell with a default value of a conditional.
\end{hbmethod}

\begin{hbtheorem}{thm:p16-zero-algorithm-soundness}{proved=annex:p16.zero-algorithm-soundness}{Termination, denominator propagation and soundness of every known return}
Both scalar extensions terminate. Every recovered kernel used by a
recursive call has the required structural meaning. Every known return
is the true requested intervention mass, and every returned whole
response kernel is identified at the input law. The source's unconditional
ancestral and additional-action identities remain valid at zero masses;
the district identities use the recovery theorem on their precise slices.
Unresolved factors in a discarded zero summand require no conditional
version or support assumption. An unresolved nonzero-weight request is
reported explicitly, rather than treated as zero.
\end{hbtheorem}

\begin{hbtheorem}{thm:p16-zero-mass-do-rules}{proved=annex:p16.zero-mass-do-rules}{The three do-calculus rules in joint-mass form}
Let \(X,Y,Z,W\) be pairwise disjoint observed sets, with respective
assignments \(x,y,z,w\), and let \(\widehat G\) be the model's
explicit latent-root DAG. Use the graph modifications of the three-rule
soundness theorem, and put
\(Z(W)=Z\setminus\operatorname{an}_{\widehat G_{\overline X}}(W)\).
Under the respective premise
\[
\begin{gathered}
 Y\mathrel{\perp_{\widehat G_{\overline X}}}Z\mid(X\cup W),
 \\
 Y\mathrel{\perp_{\widehat G_{\overline X,\underline Z}}}Z\mid(X\cup W),
 \\
 Y\mathrel{\perp_{\widehat G_{\overline X,\overline{Z(W)}}}}Z\mid(X\cup W),
\end{gathered}
\]
the following corresponding identity holds for all assignments,
including zero conditioning masses:
\[
 \begin{split}
 p_x(y,z,w)p_x(w)&=p_x(y,w)p_x(z,w),\\
 p_{x,z}(y,w)p_x(z,w)&=p_x(y,z,w)p_{x,z}(w),\\
 p_{x,z}(y,w)p_x(w)&=p_x(y,w)p_{x,z}(w).
 \end{split}
\]
These equalities do not equate arbitrary conditional versions at a null
stratum. Division recovers a conditional rule only where its stated
denominators are positive.
\end{hbtheorem}

\subsection{The completeness statement that survives}

\begin{hbdefinition}{def:p16-zero-licensed-domain}{A fully specified licensed-trace domain}
Fix either graph procedure, its topological order, and fixed choices
for reference tuples. Consider its finite scalar graph traces for
every \(x\in\mathcal A_{\mathrm{obs}}(p,X)\) and every \(y\).
For this domain definition, leave graph failures as terminal records,
and require every requested extraction table to be fully recoverable
by \(\mathsf{Recover}\); do not use the certain-action shortcut
or omit an extraction merely because another factor is zero.
The compatible laws satisfying these finite table requirements form
the \emph{licensed-trace domain} for that procedure and query.
It includes every strictly positive compatible law. It can also contain
laws with zero masses. Pointwise zero-product evaluation and the
certain-action shortcut can resolve further inputs beyond this domain.
\end{hbdefinition}

\begin{hbtheorem}{thm:p16-zero-licensed-completeness}{proved=annex:p16.zero-licensed-completeness}{Uniform completeness on licensed traces and the limit of original-action support}
Uniformly over finite nonempty alphabets and the licensed-trace domain,
the original-action supported query is identified exactly when its graph
has no hedge for the complete query. Equivalently the chosen ordinary
graph procedure succeeds. This holds for either ID or the modified Tian
procedure. A positive hedge witness proves the negative direction;
licensed recovery proves the positive direction.
This quantifies over the licensed laws; it does not classify
identification separately at each individual sparse law.

On the larger domain defined only by original-query membership in
\(\mathcal A_{\mathrm{obs}}\), or even \(\mathcal A_{\mathrm{all}}\),
the no-hedge sufficient condition and unrestricted graph-formula evaluation
are false, as the front-door counterexample proves. Therefore no
source theorem on that larger input-only interpretation is established
by merely declaring missing conditional versions arbitrary or omitting
zero observed cells. Requiring every joint action at every value instead
returns to full observed positivity by the quantified support theorem.
\end{hbtheorem}

\begin{hbtheorem}{thm:p16-zero-worked-identification}{proved=annex:p16.zero-worked-identification}{Identified zero-support examples and a necessary zero-factor cancellation}
In graph \(A\to B\to Y\) without shared inputs, let \(A\)
be fair, \(B=A\), and \(Y=B\oplus E\), with independent
\(\mathbb P(E=1)=1/4\). The only positive observed masses are
\[
 p(0,0,0)=3/8,\quad p(0,0,1)=1/8,\quad
 p(1,1,0)=1/8,\quad p(1,1,1)=3/8.
\]
The support-aware ID calculation identifies
\(P_0(Y=1)=1/4\), \(P_1(Y=1)=3/4\).
The factor for \(B\ne a\) is proved zero, so its unresolved
\(Y\) conditional is legitimately omitted. The remaining term
uses a positive conditioning prefix.

For a fully licensed nonpositive trace, take graph \(A\to Y\),
\(A\) fair and \(Y=0\). Recovery at either assigned \(a\)
returns the whole table \(Q[\{Y\}](0;a)=1\),
\(Q[\{Y\}](1;a)=0\). Finally, on the declared bow graph
\(A\to Y,A\leftrightarrow Y\), the observed law
\(p(A=0,Y=y)=1/2\) for each binary \(y\) identifies
\(P_0(Y=y)=1/2\) by the certain-action identity, although the
unmodified graph trace fails. This last conclusion is about the
specified supported action and law, not uniform identification in that graph.
\end{hbtheorem}

\begin{proofboundary}{def:p16-zero-support-scope}{Exact source interpretation and remaining completeness question}
The source's displayed action condition is preserved with both explicit
parent-stratum readings. A condition imposed on each intermediate
complementary action in its current model, at every positive parent stratum,
gives the extraction overlap check used above. The original action alone
does not imply those additional checks. The counterexample disproves
that implication; it is not an unresolved proof promise.

The finite support-aware procedure is sound wherever it returns a value,
and its graph-uniform completeness theorem has the stated licensed-trace
domain. Whether every identified mass at an arbitrary zero-support law
is resolved by these local recovery and zero-product rules is not asserted.
Nor is general do-calculus completeness on that larger law-dependent
domain inferred from the positive-law theorem. The mass-form rules and
the source's support-licensed derivations remain available, with each
division justified on its actual stratum. The fixed-binary/ternary
refinement from the earlier hedge construction is a separate question.
\end{proofboundary}
